\cleardoubleoddpage
\chapter{Background and Related Work}
\label{chap:background}

\section{Gyrokinetic Turbulence}

Gyrokinetic theory reduces the six-dimensional kinetic description by averaging over the fast
gyromotion while retaining three spatial and two velocity-space coordinates. The resulting
five-dimensional distribution function still supports nonlinear turbulence, scale interactions,
and transport. A complete derivation is outside the scope of this thesis; the multiscale framework
of \citet{abel2013multiscale} provides the physical context. For machine learning, the important
consequence is the size and anisotropy of each state: a snapshot is not a conventional two- or
three-dimensional image, and each axis has a distinct physical interpretation.

The data used here provide the distribution-function field together with a scalar flux target and
two one-dimensional spectra, denoted \code{kyspec} and \code{fluxspec}. Flux summarizes a central
transport quantity, whereas spectra expose some scale-dependent structure. They are therefore
treated as diagnostic targets rather than attempting a full reconstruction of $x_t$.

\section{Neural Surrogates for Gyrokinetics}

Recent work has begun to apply neural surrogates directly to gyrokinetic simulations. GyroSwin
extends hierarchical vision Transformers to five-dimensional fields and couples distribution and
potential representations through specialized modules \citep{paischer2025gyroswin}. Its ambition
is full-field surrogacy and physically verifiable transport estimates. GKFieldFlow instead models
three-dimensional potential fields with a spatial encoder--decoder and temporal convolutions
\citep{ashourvan2026gkfieldflow}. These systems motivate learning temporal structure but solve a
different prediction problem from the compact diagnostic-preserving latent model studied here.

The distinction matters for evaluation. A full-field model may be assessed through reconstruction,
spectral fidelity, and transport diagnostics. This thesis has no decoder and therefore makes no
full-field reconstruction claim. Its strongest evidence is an internal comparison among sequence
models that share one encoder, latent cache, split, and rollout protocol.

Neural-surrogate comparisons also require a baseline that reflects the temporal regularity of the
data. Persistence is particularly demanding when adjacent states change slowly: a learned model
must improve on copying the latest observation, not merely achieve a small absolute error. Model
choice must be separated from final assessment as well. Validation trajectories support choices
such as architecture and normalization, whereas test trajectories estimate performance only after
those choices are fixed. If representation learning has already observed a downstream test
trajectory, the complete surrogate is not held out even when the temporal model itself has not
trained on that trajectory. Model selection and performance estimation must also remain separated:
optimizing a noisy validation criterion can otherwise produce selection bias
\citep{varma2006crossvalidation,cawley2010selectionbias}. These principles motivate the trajectory-level lineage checks and the
validation-then-test protocol used in this thesis.

\section{Self-Supervised Representation Learning}

Representation learning seeks a mapping $f_\theta(x)=z$ that retains useful information while
discarding nuisance variation. Contrastive approaches often compare positive and negative pairs.
SimSiam shows that a Siamese architecture can learn without negative pairs or a momentum encoder,
provided a predictor and stop-gradient operation prevent trivial collapse
\citep{chen2021simsiam}. Given two augmented views $x^{(1)}$ and $x^{(2)}$, the shared encoder and
projection head produce $p_1$ and $p_2$, while a prediction head produces $q_1$ and $q_2$. A
symmetric negative cosine objective aligns each prediction with the stopped projection of the
other view.

PCA gives a linear projection for inspecting broad variance structure, while t-SNE constructs a
nonlinear neighbourhood-preserving visualization \citep{vandermaaten2008tsne}. Neither method is a
forecast metric: apparent clusters or gradients can depend on scaling and projection parameters.
They are therefore used only for descriptive inspection in this work.

Scientific data require conservative augmentations. Transformations that are harmless for natural
images can alter physical meaning. The present implementation uses small additive noise, amplitude
jitter, and sparse element masking. It does not use arbitrary crops or rotations. The
self-supervised objective is combined with flux and spectra supervision so the latent space is
explicitly encouraged to retain the selected diagnostics.

\section{Latent Dynamics Models}

Once snapshots are embedded, the temporal problem becomes vector sequence prediction. Persistence
copies the final context vector and provides a parameter-free baseline. An MLP-delta model predicts
a residual from the flattened context. GRUs use learned gates to update a recurrent state and are a
compact recurrent baseline \citep{cho2014gru}. Transformers replace recurrence with self-attention
over the context \citep{vaswani2017attention}. A causal mask prevents a token from accessing future
positions, and the final token representation predicts the next latent.

Autoregressive rollout introduces distribution shift: the first prediction is conditioned on true
context, whereas later predictions are conditioned partly on previous model outputs. Consequently,
one-step training loss alone does not establish rollout quality. This thesis evaluates recursive
eight-step predictions and reports each horizon as well as aggregated metrics.

\section{Scientific Machine-Learning Stack}

JAX provides accelerator-oriented array computation together with transformations for automatic
differentiation, compilation, vectorization, and device parallelism \citep{jax2026}. Flax supplies
the neural-network module abstraction used here, and Optax supplies composable optimization
updates \citep{flax2026,optax2026}. The implementation uses pure JAX computation at the model and
training boundary, while Python, NumPy, HDF5, and KvikIO-compatible readers handle file I/O. This
division keeps I/O outside compiled functions and permits CPU smoke tests alongside GPU training.
